\section{Análise Comparativa Global e Benchmarking}

Nesta seção, apresentamos os resultados consolidados da avaliação final dos modelos no conjunto de teste (\textit{hold-out}), que representa dados inéditos para as arquiteturas. O objetivo principal desta análise é identificar a arquitetura com maior capacidade de generalização e resiliência ao desbalanceamento de classes, utilizando o Coeficiente de Correlação de Matthews (MCC) como métrica primária de seleção, conforme estabelecido no protocolo experimental.

Os resultados detalhados para as principais variações de datasets (V1, V9 e V10) e arquiteturas estão apresentados na Tabela \ref{tab:results_full_metrics}. Observa-se que o modelo \textbf{DenseNet} na versão de dataset V10 obteve o melhor desempenho global, alcançando um MCC de 0,6461 e uma acurácia de teste de 69,58\%. Este resultado é significativamente superior aos modelos de base inicial (V1), demonstrando que o ajuste fino de hiperparâmetros via otimização bayesiana e o tratamento de iluminação na V10 foram eficazes.

\begin{table}[!h]
    \centering
    \begin{threeparttable}
    \caption{Métricas consolidadas de desempenho e complexidade no conjunto de teste.}
    \label{tab:results_full_metrics}
    \begin{tabular}{lllll}
\toprule
Model & MCC & F1-Macro & Kappa & GFLOPs \\
\midrule
CONVNEXT & \textbf{0,5503} & \textbf{0,6194} & \textbf{0,5426} & 4,4702 \\
DENSENET & 0,5501 & 0,6050 & 0,5390 & 2,8651 \\
EFFICIENTNET & 0,5159 & 0,5924 & 0,5148 & \textbf{0,4010} \\
VGG & 0,4917 & 0,5474 & 0,4882 & 19,5446 \\
SWIN & 0,3931 & 0,4800 & 0,3889 & 4,5090 \\
VIT & 0,3910 & 0,4721 & 0,3902 & 16,8669 \\
RESNET & 0,3238 & 0,4250 & 0,3187 & 3,6711 \\
INCEPTION & 0,2484 & 0,3435 & 0,2467 & 2,8473 \\
\bottomrule
\end{tabular}

    \begin{tablenotes}
            \small 
            \item[] Fonte: Do autor.
        \end{tablenotes}
        \end{threeparttable}
\end{table}

Um achado relevante é o desempenho das arquiteturas modernas baseadas em Transformers (\textit{Vision Transformers} - ViT e Swin Transformer). Apesar de possuírem uma complexidade paramétrica elevada (86,0M para o ViT), esses modelos apresentaram MCC inferior (máximo de 0,4328) em comparação com as CNNs tradicionais como DenseNet e ConvNext. Este comportamento corrobora a literatura de que, para conjuntos de dados de domínio específico e tamanho moderado, o viés indutivo (\textit{inductive bias}) de localidade e invariância de translação das CNNs ainda oferece uma vantagem competitiva sobre a flexibilidade global dos mecanismos de atenção dos Transformers.

Outro ponto crítico identificado foi o fracasso da versão experimental V8, que utilizou uma expansão para 15 canais espectrais via filtros de Gabor. Conforme observado na Tabela \ref{tab:results_full_metrics}, todos os modelos apresentaram uma queda drástica de performance na V8, com MCCs próximos de zero. Este fenômeno sugere que a inserção de \textit{features} artesanais ruidosas ou a alta dimensionalidade sem o volume de dados correspondente prejudicou o aprendizado das representações latentes pelas redes, validando a abordagem de \textit{end-to-end learning} sobre os pixels RGB brutos.
